% Financial Mathematics
% Finance Final for Spring 2026
%----------------------------------------------
\documentclass[12pt]{exam}

\usepackage{setspace}
\setstretch{1.15}
\nopointsinmargin
\pointformat{}

\usepackage{problemset}

%\usepackage[margin=1in]{geometry}

% Uncomment ONE of these:
\noprintanswers
%\printanswers

\begin{document}

\FontEasyRead

\Heading
{Financial Literacy}
{Final Exam}
{60 Minutes, 100 points plus potential 5 bonus points.
  Up to 4 handwritten pages (2 pieces of paper) of notes and a calculator are permitted.
  Answer all multiple-choice questions by selecting the best answer.
  Answer all short-response questions clearly, using correct spelling, punctuation, and grammar.
  Show all work for financial math questions, including any pertinent timeline sketches.
Don't leave any blank answers.}

\section*{Part I: Multiple Choice (2 points each)}

\begin{questions}

  \begin{mcquestion}[2]
    Buying an expensive brand-name product just because others are buying it is an example of:
    \begin{choices}
      \choice Loss aversion
      \choice Endowment effect
      \CorrectChoice Herd behavior
      \choice Budget constraint
    \end{choices}
  \end{mcquestion}

  \begin{mcquestion}[2]
    When a store advertises a product as “normally \$100, now \$60,” the phony "normal" price is
    primarily used to:
    \begin{choices}
      \choice Reflect production cost
      \choice Signal quality
      \CorrectChoice Anchor the consumer’s perception of value
      \choice Increase supply
    \end{choices}
  \end{mcquestion}

  \begin{mcquestion}[2]
    A benefit such as health insurance provided by an employer is best described as:
    \begin{choices}
      \choice Gross income
      \choice Net pay
      \CorrectChoice Non-wage compensation
      \choice Payroll tax
    \end{choices}
  \end{mcquestion}

  \begin{mcquestion}[2]
    Which of the following best describes a progressive tax system like the
    current U.S. Federal Income Tax Code?
    \begin{choices}
      \choice Everyone pays the same dollar amount
      \choice Everyone pays the same percentage of income
      \CorrectChoice Higher incomes are taxed at higher rates
      \choice Lower incomes are taxed at higher rates
    \end{choices}
  \end{mcquestion}

  \begin{mcquestion}[2]
    Which type of account typically offers the easiest access to funds for daily transactions?
    \begin{choices}
      \CorrectChoice Checking account
      \choice Certificate of deposit (CD)
      \choice Retirement account like a 401(k) or 403(b)
      \choice Stock brokerage account
    \end{choices}
  \end{mcquestion}

  \begin{mcquestion}[2]
    Which budgeting strategy assigns to every dollar of income a specific purpose?
    \begin{choices}
      \choice 50/30/20 rule
      \choice Pay-yourself-first method
      \choice Incremental budgeting
      \CorrectChoice Zero-based budgeting
    \end{choices}
  \end{mcquestion}

  \begin{mcquestion}[2]
    Which retirement account provides \underline{tax-free} withdrawals during retirement?
    \begin{choices}
      \choice Traditional IRA
      \choice 401(k)
      \CorrectChoice Roth IRA
      \choice Pension plan
    \end{choices}
  \end{mcquestion}

  \begin{mcquestion}[2]
    Diversification helps investors primarily by:
    \begin{choices}
      \CorrectChoice Reducing volatility by spreading investments
      \choice Increasing guaranteed returns
      \choice Eliminating losses
      \choice Avoiding taxes
    \end{choices}
  \end{mcquestion}

  \begin{mcquestion}[2]
    The amount a policyholder must pay out-of-pocket before insurance coverage begins is called:
    \begin{choices}
      \choice Premium
      \choice Copayment
      \CorrectChoice Deductible
      \choice Claim
    \end{choices}
  \end{mcquestion}

  \begin{mcquestion}[2]
    Liability insurance primarily protects against:
    \begin{choices}
      \choice Damage to your own property
      \choice Medical expenses only
      \CorrectChoice Financial loss from harming others or their property
      \choice Investment losses
    \end{choices}
  \end{mcquestion}

  \begin{mcquestion}[2]
    Carrying a balance on a credit card typically results in:
    \begin{choices}
      \choice Higher credit score automatically
      \choice No additional cost
      \CorrectChoice Interest charges on the unpaid balance
      \choice Lower minimum payments
    \end{choices}
  \end{mcquestion}

  \begin{mcquestion}[2]
    Which benefit is commonly offered by employers in addition to salary?
    \begin{choices}
      \choice APR
      \choice Credit score increase
      \choice Sales tax refund
      \CorrectChoice Life and disability insurance
    \end{choices}
  \end{mcquestion}

  \begin{mcquestion}[2]
    Which of the following is an example of a variable expense?
    \begin{choices}
      \CorrectChoice Grocery bills
      \choice Monthly rent
      \choice Car insurance premium
      \choice Student loan payment
    \end{choices}
  \end{mcquestion}

  \begin{mcquestion}[2]
    Which loan typically has the shortest repayment period (term)?
    \begin{choices}
      \choice Mortgage loan
      \choice Student loan
      \choice Auto loan
      \CorrectChoice Payday loan
    \end{choices}
  \end{mcquestion}

  \begin{mcquestion}[2]
    What is the primary purpose of a credit report?
    \begin{choices}
      \choice To calculate income taxes
      \CorrectChoice To show a person's borrowing and repayment history
      \choice To track investment returns
      \choice To determine insurance premiums only
    \end{choices}
  \end{mcquestion}

  %%%%%%%%%%%%%%%%%%%%%%%%%%%%%%%%%%%%%%%%%%%%%%%%%%%%%%
  %  \newpage
  \section*{Part II: Short Answer (10 points each)}

  \begin{shortanswer}[10]{2in}
    Morgan uses a credit card but only makes minimum payments each month, and occasionally
    forgets to make any payment at all. Explain:
    \begin{itemize}
      \item How interest and penalty fees affect the total cost of purchases
      \item The likely impact on Morgan’s credit score resulting from his behavior
      \item A better strategy for using credit responsibly
    \end{itemize}
  \end{shortanswer}
  \begin{solution}
    Interest increases total repayment over time. Carrying balances may increase utilization and hurt credit score. Better strategies include paying in full and on time.
  \end{solution}

  \begin{shortanswer}[10]{1.5in}
    Leo is trying to choose between a high-deductible (say \$5,000) and low-deductible
    (say \$250) insurance plan.  Explain:
    \begin{itemize}
      \item The trade-off between premium level and deductible size
      \item When a high-deductible plan might make sense
      \item When a low-deductible plan might be more appropriate
    \end{itemize}
  \end{shortanswer}
  \begin{solution}
    High deductibles usually mean lower premiums and vice versa.
    High-deductible plans work for healthy individuals with savings.
    Risk is high out-of-pocket costs in
    emergencies.
  \end{solution}

  \begin{shortanswer}[10]{2in}
    Explain the concept of an investor's ``time horizon'' and how it impacts his investment
    strategy decisions.  What kinds of investments are best suited for someone who is just
    entering the workforce?  For someone who is just about to retire?
  \end{shortanswer}
  \begin{solution}
    How long they can keep their money invested before needing to access it.  Longer horizons can
    take more risk for higher returns (e.g. stocks).  Shorter horizons prioritize safety and
    liquidity (e.g. bonds, cash).
  \end{solution}

  \begin{shortanswer}[10]{2in}
    Someone is asking for your advice about payment methods.  Explain the general characteristics, and pros
    and cons, of these payment methods:
    \begin{itemize}
      \item Debit card
      \item Credit card
      \item Cash
      \item Mobile payment apps (e.g. Apple Pay, Venmo, etc.)
    \end{itemize}
  \end{shortanswer}
  \begin{solution}
    Debit cards draw directly from checking accounts, no interest but risk of overdraft. Credit
    cards offer convenience and rewards but can lead to debt if not paid in full. Cash is widely
    accepted and helps with budgeting but can be lost or stolen. Mobile apps are convenient for
    peer-to-peer payments but may have security risks.
  \end{solution}

  %%%%%%%%%%%%%%%%%%%%%%%%%%%%%%%%%%%%%%%%%%%%%%%%%%%%%%%%%%%%%%
  \section*{Part III: Financial Math (10 points each)}

  \begin{problem}[10]{2in}
    You purchase \$3,000 worth of a particular stock.  The first year, the price of the
    stock rises 6.5\%.  The next year, the price \underline{falls} 8.3\%.  The third year the price
    rises by 11.2\%.

    What is the value of your stock holding after three years? What was
    your \underline{average annual} rate of return over the three year period?
  \end{problem}
  \begin{solution}
    \begin{align*}
      \text{Year 1:} \quad & 3000 \times 1.065 = 3195 \\
      \text{Year 2:} \quad & 3195 \times 0.917 = 2930.72 \\
      \text{Year 3:} \quad & 2930.72 \times 1.112 = 3258.24
    \end{align*}

    \vspace{0.5\baselineskip}

    \text{Average annual return:}
    \begin{align*}
      \left(\frac{3258.24}{3000}\right)^{1/3} - 1 &\approx 0.0287 \\
      &= 2.87\%
    \end{align*}
  \end{solution}

  \begin{problem}[10]{2in}
    You see a \$1,000 zero coupon bond that matures in 6 months.  You buy it at the market
    price of \$963.

    When you redeem the bond at maturity for \$1,000, what rate of interest rate
    did you earn over the 6 month period?  What would that rate be if you annualized it?
  \end{problem}
  \begin{solution}
    \begin{align*}
      \text{6-month return:} \quad & \frac{1000 - 963}{963} \approx 0.0384 \\
      &= 3.8\%
    \end{align*}

    \vspace{0.5\baselineskip}

    \text{Annualized return:}
    \begin{align*}
      (1 + 0.0384)^2 - 1 &\approx 0.0783 \\
      &= 7.8\%
    \end{align*}
  \end{solution}

  \begin{problem}[10]{2in}
    Sam and Dave both invest \$1{,}000 at 5\% annual interest, compounded annually.
    Sam leaves the money invested for 3 years.
    Dave withdraws all interest at the end of \underline{each year} instead of leaving it in the account.

    Calculate:
    \begin{itemize}
      \item Sam’s final balance after 3 years, and how much interest he's earned in total
      \item The total amount of interest Dave earns over 3 years
    \end{itemize}
    Who ends up with more money, and why?
  \end{problem}
  \begin{solution}
    Sam:
    Year 1: \$1{,}000 $\times$ 1.05 = \$1{,}050
    Year 2: \$1{,}050 $\times$ 1.05 = \$1{,}102.50
    Year 3: \$1{,}102.50 $\times$ 1.05 = \$1{,}157.63

    Dave:
    Earns \$50 each year (5\% of \$1{,}000), so total interest = \$150.

    Sam ends with \$1{,}157.63, while Dave has \$1{,}150 total.

    Sam ends up with more because of compounding—earning interest on interest.
  \end{solution}

  %%%%%%%%%%%%%%%%%%%%%%%%%%%%%%%%%%%%%%%%%%%%%%%%%%%%%%
  \section*{Bonus Question - 5 points}

  \begin{shortanswer}[5]{3in}
    What is one financial lesson you've learned from this course that you think will be most
    useful in your life?
  \end{shortanswer}
  \begin{solution}
    Neither a borrower nor a lender be.

    For loan oft loses both itself and friend,

    And borrowing dulls the edge of husbandry.
  \end{solution}

\end{questions}

\end{document}
